\section{Merge Sort}
O MergeSort é um algoritmo de ordenação que segue o paradigma de divisão e conquista. A ideia principal é dividir o vetor original em dois vetores menores, ordená-los recursivamente e, em seguida, intercalá-los para formar o vetor final ordenado.

Um dos passos fundamentais para a implementação do MergeSort é a intercalação de dois vetores ordenados, processo para o qual usar-se-á o termo em inglês Merge. O problema consiste em combinar os elementos desses vetores em um único vetor também ordenado. Como ambos os vetores já estão organizados, é possível explorar essa característica para criar um algoritmo eficiente, que percorre os elementos de forma sequencial, comparando-os e adicionando o menor ao vetor resultante.

\lstinputlisting[language=C]{funcoes-c/g_merge.txt}
% \begin{algorithm}
%     \caption{Merge}\label{alg-merge:cap}
%     \DontPrintSemicolon
%     \Entrada{$inicio$, $meio$, $fim$ e $v$, com $v[inicio..meio-1]$ e $v[meio..fim-1]$ ordenados}
%     \Resultado{o vetor $v[inicio..fim-1]$ em ordem crescente}
%     \Inicio{
%         $aux \gets CriaVetor(fim - inicio)$ \tcp*{Cria um vetor de tamanho $fim - inicio$}
%         $i \gets inicio$\;
%         $j \gets meio$\;
%         $k \gets 0$\;
%         \Enqto{$i < meio$ \textbf{ou} $j < fim$}{
%             \eSe{$j \geq fim$ \textbf{ou} $(i < meio$ \textbf{e} $v[i] \leq v[j])$}{
%                 $aux[k] \gets v[i]$\;
%                 $i \gets i + 1$\;
%             }{
%                 $aux[k] \gets v[j]$\;
%                 $j \gets j + 1$\;
%             }
%             $k \gets k + 1$\;
%         }
%         \Para{$i = inicio$ \Ate $\,fim - 1\,$}{
%             $v[i] \gets aux[i - inicio]$\;
%         }
%     }
% \end{algorithm}

Agora o algoritmo Merge será usado como sub-rotina do algoritmo principal.

\lstinputlisting[language=C]{funcoes-c/g_merge_sort_recursivo.txt}
\lstinputlisting[language=C]{funcoes-c/g_merge_sort.txt}
% \begin{algorithm}
%     \caption{MergeSort}\label{alg-merge-sort:cap}
%     \DontPrintSemicolon
%     \Entrada{um intervalo [$inicio$, $fim$) e um vetor $v$, com $0 \leq inicio,\, fim \leq |v|$}
%     \Resultado{o vetor $v[inicio..fim-1]$ em ordem crescente}
%     \Inicio{
%         \Se{$inicio < fim - 1$}{
%             $meio \gets (inicio + fim) / 2$\;
%             MergeSort($inicio$, $meio$, $v$)\;
%             MergeSort($meio$, $fim$, $v$)\;
%             Merge($inicio$, $meio$, $fim$, $v$)\;
%         }
%     }
% \end{algorithm}

Observe que para ordenar um vetor $v$ de tamanho $n$, deve-se executar MergeSort($0$, $n$, $v$).

\subsection{Prova de correção}
Seja $n$ o tamanho do vetor $v$. Vamos provar que o algoritmo é correto por indução em $n$. Se $n = 1$ a condição da instrução \textbf{se} é falsa e nenhuma movimentação é realizada no vetor. Como o vetor já estava ordenado por definição, então o algoritmo é correto para este caso.

Agora suponha que $n > 1$ e que o algoritmo funciona para qualquer vetor de tamanho menor que $n$. Neste caso, o algoritmo divide o vetor em duas partes de tamanho menor que $n$:
\[
  \overbrace{v[0..\lfloor n/2\rfloor-1]}^{\lfloor n/2\rfloor} \hspace{0.5cm}\text{ e }\hspace{0.5cm} \overbrace{v[\lfloor n/2\rfloor..n-1]}^{n - \lfloor n/2\rfloor}
\]
Por hipótese, podemos supor que o algoritmo vai ordenar corretamente as duas partes. Por fim, o algoritmo Merge vai intercalar corretamente as duas partes, garantindo que o vetor completo fique em ordem.

Portanto, pelo princípio da indução, o algoritmo MergeSort é correto para vetores de qualquer tamanho.

\subsection{Complexidade de tempo e espaço}
Suponha, sem perca de generalidade, que a execução do algoritmo inicia com um vetor de tamanho $n = 2^k$, com $k \in \mathbb{N}$. A cada chamada recursiva do algoritmo, o tamanho do intervalo diminui exatamente pela metade. Ou seja, na primeira chamada o intervalo terá tamanho $2^{k-1}$, na segunda $2^{k-2}$, e assim por diante. Sendo assim, a arvore de recurção terá exatamente $k+1$ níveis, como ilustra a tabela \ref{tab:merge-sort-interval-lenght}.

\begin{table}[!h]
    \centering
    \Caption{\label{tab:merge-sort-interval-lenght}Tamanho do intervalo por nível da árvore de recursão - MergeSort}
    \begin{tabular}{c|c|c|c|c|c|c}
        Nível & $1$ & $2$ & $3$ & $\cdots$ & $k$ & $k+1$ \\
        \hline
        Tamanho do intervalo & $2^k$ & $2^{k-1}$ & $2^{k-2}$ & $\cdots$ & $2$ & $1$ \\
    \end{tabular}
\end{table}

No i-ésimo nível da árvore de recursão, temos $2^{i-1}$ intervalos de tamanho $2^{k-i+1}$. O tempo gasto em cada um desses intervalos é o tempo gasto na operação de Merge que, por sua vez, é linear no tamanho do intervalo. Mas a soma dos tamanhos de todos os intervalos de um mesmo nível é igual à $n$:
\[
    2^{i-1} \cdot 2^{k-i+1} = n
\]
Logo, o tempo total gasto em cada nível é proporcional a $n$. Como existem $k+1$ níveis, então o tempo de execução do algoritmo é $O(n\cdot k)$. Por fim, para deixar tudo em função de $n$, observe que:
\[
    n = 2^k \Leftrightarrow k = log_2(n)
\]
Portanto, a complexidade temporal do MergeSort é $O(n\cdot log_2(n))$.